
The emergence of the internet enabled people all over the globe to connect with each other in ways unseen before. Through social networks, information exchange, discussions, and opinion expression became easier, more widely spread, and increasingly instantaneous. Social networks and their associated platforms also play an increasingly important role in financial markets \citep{Deng2024, Gopal2024,Bartov2023,Peng2024,Gu2020}. \str{Especially Twitter received notable attention \citep{lehrer2021does, Gu2020,Behrendt2018, Yang2015}, but also StockTwits \citep{Audrino2020, chen2022sonic} and sentiment from across sources \parencite{audrino2024does, Gan2020, Jiao2020}. }
However, the impact of Reddit on society and markets is less frequently studied even though the social network is the 5th most visited social media network globally as of March 2025 according to \href{https://www.similarweb.com/}{similarweb.com}. In the USA, home to the largest financial market, Reddit.com is even the 2nd most visited social media network, being more popular than LinkedIn, Tiktok and X. 
Latest since the GameStop (GME) market frenzy in January 2021 \footnote{GameStop, a brick and mortar video game retailer, struggled with competition from digital distribution services and the negative economic impact of the COVID-19 pandemic. As a result, the stock price declined, and many hedge funds short-sold the stock. At that time, WSB users began to buy and hold GME stocks in bulk. Excess demand combined with limited availability drove the stock price and ultimately triggered a short squeeze. The stock price of GME rose from \$ 0.70 on 3 April 2020 to \$18.80 on 31 December 2020 and further to \$483 on 28 January 2021. Short sellers like Melvin Capital suffered considerable losses in this short squeeze.} , which was driven by users of the subreddit `Wallstreetbets' (WSB, \hyperlink{https://www.reddit.com/r/wallstreetbets/}{reddit.com/r/wallstreetbets}), the platform received notable attention in society and academia \citep{Pedersen2022,Klein2022,Umar2021,molders2025understanding}. \str{While attention within a community is well known to impact stock prices \citep{Andrei2014}}, the frenzy showed that users of social networks have real impact on financial markets. The effect of the market frenzy was so detrimental that even a US senate hearing took place, during which they questioned the main actors involved in the market frenzy, including one of the Reddit users hyping GME. In the aftermath of the GME frenzy, it became obvious that the `internet crowd' represents a new force driving financial markets. Constantly new stocks enter the focus of the `internet crowd'\str{, leading regularly to herding behaviour and market excesses, impacting asset prices \citep{Aloosh2021, Costola2021}}. These assets range from GME, Virgin Galactic, and Palantir, to Gitlab, AMC and BB, see for example \hyperlink{https://www.fool.com/investing/stock-market/types-of-stocks/meme-stocks/}{Motley Fool}, \hyperlink{https://www.benzinga.com/money/best-meme-stocks}{Benzinga}, or \hyperlink{https://finance.yahoo.com/news/12-best-meme-stocks-buy-191420607.html}{Yahoo Finance}. 
Due to its continued relevance, the U.S. Securities and Exchange Commission (SEC) introduced a new term for stocks being on the radar of the internet crowd: the `meme stock phenomenon' \citep{Securities2021}. 

The DeGroot type model of \citet{Pedersen2022} provides a theoretical framework for understanding (meme) stock prices discussed and traded by users of social networks. The model provides a theoretical framework for network users \str{behaviour}. It being a DeGroot model implies that users of the network change their \str{behaviour in relation to} a stock as a result of information they receive from the network, which translates then into price changes. However, this theoretical model is not yet thoroughly backed up by empirical evidence. It raises questions like what is driving network users' posting behaviour, which users are driving the network communication, how does this spill over into the financial markets, and how large is the effect of social network users on the actual financial market?

In this paper, we show that the change in behaviour of individual users on a social network explains market dynamics, thereby providing the empirical evidence to support the theoretical model by \citet{Pedersen2022}. For the study we hand collect 3.8 million posts from 502,650 users from WallStreetBets, a subreddit. To analyse this large scale network dataset, we propose the Interpretable Social Network and Market Behaviour modelling approach (ISMB), which combines logistic regression with network variables under sparsity and change point detection. %Part of the model is the construction of change point based variables to quantify the change of behaviour of users and study the relationship with market movements. 
The modelling approach captures heterogeneous individual user behaviour and links this through innovative change point modelling to the stock market. It is computationally scalable and accounts both for influential nodes as well as followers. The approach allows to study the individual relationship between any two nodes and explains their behaviour via the network variables embedded in it. This allows us to determine what drives network users' posting behaviour and to study their change in behaviour over time. We find that the change in behaviour of users on the social network aids the explanation of market movements in various ways: 1) We show that embedding the change in behaviour of users on the social network as an exogenous variable into an autoregressive model, explains a larger portion of the market dynamics than the ones without. %For return dynamics, the adjusted $R^2$ of said models is 630\% larger. 
In combination with the theory outlined in \citet{Pedersen2022}, this suggests that retail investors impact meme stock prices with their change in behaviour. 2) We find that influential users' changes in behaviour result after 6 hours in a significant effect upon market dynamics. This makes sense since a post by a single influential person first needs to reach a credible number of investors to lead too market reactions. While it is extremely challenging to detect change points in large scale networks quickly \citep{zwetsloot2025online,Yu2022monitoring}, this suggests that influential users change in behaviour could act as an early warning sign for meme stock market reactions, if identified early enough. We also find that several influential users exhibit a change in behaviour 1-2 months before the market frenzy. 3) Lesser influential users change in behaviour, if occurring in bulk, ripples instantly and with a lasting effect of various hours into the financial markets. Their change in behaviour explains the largest portion of market dynamics, even when controlling for network activity. This finding paves the way towards the construction of variables for stock market movements based on social network users behaviour. 

%The GameStop (GME) market frenzy of 2021 brought an unseen situation to the financial markets. An internet community hyped up the stock of a company and caused its price to skyrocket whereas the companies fundamentals did not warrant such a development. At the same time hedge funds were betting against a price increase of GME via put options. The combination of both caused a short squeeze in the market which led to the bankruptcy of a hedge fund and a David-against-Goliath hype on the Wallstreetbets forum on Reddit leading to the suspension of GME trading on the popular Robinhood trading app. The saga peaked with an investigation by the United States Congress into alleged market manipulation by users of Wallstreetbets and misbehaviour of Robinhood representatives. 

%Indeed attention within a community is well known to impact stock prices \citep{Andrei2014}. Lately the effect of users on stock prices of various kind is so strong and widespread, that it led the U.S. Securities and Exchange Commission (SEC) to introduce a new term for stocks being on the radar of the internet crowd: the `meme stock phenomenon' \citep{Securities2021}. The hype surrounding `meme stocks' leads regularly to herding behaviour and market excesses, impacting asset prices \citep{Aloosh2021, Costola2021}. The ongoing effect from Reddit on financial markets via meme stocks drive the need to analyse the behaviour of users of Reddit forums and quantify their impact on the market. 

\str{Our proposed ISMB modelling approach was specifically designed for the modelling of individual users behaviour in large scale social networks while ensuring interpretability. Major challenges we tackle are the estimation and scalability, identification of changes in behaviour, as well as the large scale of the network with heterogeneous communications (links). }
\str{Link formation in social networks is a widely studied topic. Random network models such as Exponential Random Graph Models and Latent Variable Models have been widely studied for such tasks \citep{kim2018review,Deng2024}. Not all of these models are applicable for our purpose, as, for example, many become intractable, which makes them (only) suitable for small networks \citep{Yon2021}. Recently, there have been studies using deep learning models for link prediction \citep{barros2021survey,yu2023towards,beladev2023graphert,hui2024domain}. However, these models are not easily interpretable, and therefore provide little insight into the drivers of individuals posting behaviour. For large networks another option is to use logistic regression-based dynamic network models \citep{agaskar2013alarm,almquist2013dynamic}, because these provide a scalable framework for inference on large-scale temporal networks \citep{mallik2019stable}. In addition, these models naturally allow for incorporation of network dynamics through the autoregressive modelling structures. Furthermore, \textcite{bhattacharya2019coevolution} found that variables such as the past posting frequency and a nodes' out-degree matter for their future communication behaviour, opening the pathway to predict the probability of communication between two users based on their past behaviour and relevance to the network. This motivates our adoption of an individual-level logistic regression type model which is scalable as well as interpretable and allows naturally for incorporation of network dynamics.} 

\str{The heterogeneous nature of nodes implies that their communication probabilities with other users may also be influenced by different sets of determinants. Therefore, not all parameters included in the model may be relevant for every users' behaviour, hence the system may experience sparsity. In fact, many networks are characterised by huge dimensionality, but have sparse links \citep{paula_econometrics_2017}. The high-dimensional nature of networks challenges the joint analysis of nodes and identification of important links. Methods from sparse modelling, originally developed for linear models \citep{tibshirani_regression_1996, fan_variable_2001, zou_adaptive_2006}, were increasingly used for network analysis to handle the sparse nature of the underlying models \citep{Barigozzi2019, Trimborn2022a, Trimborn2022, Zhang2023}. The issue of sparsity within social networks concerns various modelling challenges, ranging from link prediction \citep{Trimborn2022} to change point detection \citep{chen2024monitoring}. Likewise, the ISMB modelling approach incorporates sparse estimation to account for the observed sparse structure. }

\str{Link formation in social networks is strongly interlinked with the role users have within networks. Naturally users `trust' certain other users more, which can lead to the emergence of opinion leaders \citep{Li2011}, also termed influencers \citep{han2020importance, chen2022sonic}. Influencers are typically characterized by their topological position in the network \str{\citep{Aral2014}}, for example as measured by various centrality measures \citep{gelper2021competition}. The difference in user relevance on the social network is known to explain user behaviour on a social network \citep{Lu2013,zhang2018direct} as well as its own behaviour and the one of its peers. It has also been shown that influencers do influence one another \citep{Ji2025}. Our study builds on these findings by modelling separate models for vital and non-vital users in the network (one may think of this terminology as influencers and non-influencers), enabling us to study the effect of these different roles in the network on the markets.}

\str{We place particular emphasis on changes in the social network and propose a greedy change point algorithm, as in \citet{bai2023multiple}, to identify changes in each link specific model. Changes are typically identified by comparing models before and after a possible change point. The overarching aim of change point detection is to identify a point in time where the network or a node in the network exhibits a change in behaviour \citep{ghezzi2024fast}. Within the field of network analysis, there is an extensive literature \citep{kendrick2018change, woodall2017overview}.}  
\str{Most of this literature considers networks with a small number of nodes, typically no more than a few thousand \citep{chen2024monitoring,Yu2022monitoring}. The huge size of social networks in combination with their typical sparsity often challenges the existing algorithms for change point detection due to their computational complexity. In addition, most literature focuses on so-called global change point detection, identifying time points where the whole network, or a substantial part of the network, changes. An example is \citet{zwetsloot2025online} who identify change points in the same GME data, however, these are at a network community level, providing no-change identification at the nodal level. It is well known that nodes have heterogeneous effects on the network \citep{chen2023community}, especially highly influential nodes, hence the target of this study requires change point detection on the nodal level. Also, as networks are highly dynamic, users may change their behaviour multiple times within our observation period, hence we need algorithms for multiple change point detection \citep{niu2016multiple,bai2023multiple}. This motivates our adoption of multiple change points detection at the node level.} 

%Our proposed ISMB modelling approach is a scalable approach based on individual users contrary to a lot of empirical work which uses aggregated network metrics such as network activity %, the aggregated number of posts at a given time, 
%have been shown to be important factors 
%for modelling the influence of social networks on financial markets \citep{Audrino2020,Gan2020,Jiao2020}. However, aggregated metrics mask interactions on the individual node level, whereas network theory shows that these interactions contain relevant information to explain the users actions \citep{Pedersen2022}. This theory builds up on the rational that people in a network update their beliefs by listening to other people in their social network \citep{DeGroot1974}. The DeGroot model states that users share their opinion with a certain probability with others, resulting in the updating of their opinions. The co-occurrence of changes in beliefs and behaviour in a social network has been observed by \textcite{Bailey2024}. This implies a need to model the effects of social networks on markets at the individual node-level rather than with aggregated measures. Drawing on this, we develop a node-specific model of posting behaviour that enables us to study what drives individuals in the network to change their behaviour. Building on the theory of influencers and investors who follow them \citep{Pedersen2022}, we build a model to investigate the relationship between individuals from these two groups of nodes, which we call Vital Nodes (VN) and Non-Vital Nodes (NVN) in the network. This also aligns with \textcite{Lu2013}'s findings that in addition to a node's past behaviour, the nodes' relevance for the network matters as well.
%Node level analysis allows to study user posting behaviour on Reddit, how they influence each other, and what drives changes in their behaviour over time. Understanding individual users behaviour enables one to identify users vital for the network, study their impact on other users, study and understand what drives the overall changes in the network, which, as it is the case for meme stocks, influences the stock markets.


%Estimation of individual-level network activity models is non-trivial. Networks generally consists of a huge number of nodes (502,650 for the network we analyse on WSB) which typically show highly dynamic and heterogeneous posting behaviours. Exponential Random Graph Models have been widely studied for link prediction \citep[e.g.][]{Deng2024}, however, for large networks they become intractable, which makes them particularly suitable for small networks \citep{Yon2021}. This motivates our adoption of a link-specific model, containing variables representing the involved user's own past behaviour and the ones of its peers, allowing for an interpretable model of communication behaviour, capturing nodes heterogeneity. This also enables us to model individual-level changes, aligning our empirical model with the theoretical framework of \citet{DeGroot1974} and \citet{Pedersen2022}. %\hl{As an added benefit, you can implement link-specific models in parallel, ensuring scalability.} %We use sparse modelling methods \citep{Barigozzi2019, Trimborn2022a, Trimborn2022, Zhang2023} to identify the relevant determinants of communication between two nodes in a data-driven manner. We include variables representing the node and VNs' behaviour, such as posting frequency, out-degree, and the nodes' past behaviour, in line with findings of \citet{bhattacharya2019coevolution, Lu2013}.

%We place particular emphasis on changes in the social network and propose a greedy change point algorithm, as in \citet{bai2023multiple}, to identify changes in each link specific model. We use these changes to study how network dynamics explain market movements. Although prior studies document that changes on the individual level are important (see \citet{Pedersen2022,Aral2014,Li2011,bhattacharya2019coevolution}, among others) it is not clear how they can be linked to the stock market empirically. We develop change point variables based on the changes identified in each individual model. With these change point variables, we link the individual-level network behaviour to the stock market and show that users change in behaviour explains market movements. %The aggregated variable measures the overall networks user’s belief changes. Hopefully capturing the users updating of thei opinions and the ascoiatcate change in behvioar (de GROOT)

%These empirical necessities led us to the construction of the Interpretable Social Network and Market Behaviour Model (ISBM). It consists of three building blocks, capturing the heterogeneous individual users' behaviour and linking it through innovative change point modelling to the stock market. %We design this model given existing methods often assume general network structures. However, each online social networks is unique in its design and its structure is determined by its architecture, functionalities, and intended use. 
%We tailor the ISBM so that it is specifically designed for analysing user behaviour on Reddit and link the results to market changes. By this, we contribute a model to analyse meme stock market movements and their link with the subreddits the meme stock is discussed in. The model developed in this study accounts for two types of users, vital and non-vital. By this, we take into account that users posting behaviour is naturally driven by different factors subject to their role in the network. \emph{In the models first block}, we build link prediction models for the individual user pairs, distinguishing between Vital and Non-Vital Nodes. The model predicts the probability for a user to post in the next period subject to their own past posts, in- and out-degree, overall network activity, responses (activity) to their past posts as well as the hour of the day. In the case of non-vital users we add parameters to take into account that they may pay considerable attention to certain vital users (follow them) which may impact the probability of them posting. The model is flexible and allows for including other metrics if relevant. When the model is applied to a subreddit focused upon asset discussions such as `Wallstreetbets', the model is extended with a parameter accounting for the opening hours of the market. The model uses LASSO to ensure model interpretability. %Focusing on individual users implies that a tailored model per user is needed to understand the drivers behind their behaviour. Not all parameters included in the model may be relevant for every user's behaviour, hence the system may experience sparsity. We observe that the posting frequency on Reddit is sparse, see section \ref{sec:data_desc}, which motivates the use of a sparse modelling framework for the user network. 
%It is well known that social networks and user behaviour are highly dynamic over time. The ISBM model has to account for this network property as well. Focusing on the subreddit `Wallstreetbets', we observe that it went through various stages of activity over time, see section \ref{sec:data_desc}. The discussions about GME increased in frequency from summer 2020 and peaked during the market frenzy in January 2021 after which the discussions about GME became less frequent. This dynamic network behaviour motivates the existence of change points on Reddit. Consequently, %we present a changepoint detection framework for the model we are introducing, based upon a likelihood based search algorithm as used in \textcite{bai2021multiple}. 
%\emph{in the second block of the ISBM}, we are concerned with the determination of change points in the communication between users. To identify if and when users change their behaviour, we specify a change point detection framework based upon a greedy algorithm, similar to \citet{bai2023multiple}. 

%\emph{In the third block of the ISBM}, we investigate how the detected changes in user behaviour link to market movements. We regress the obtained change points in the behaviour of Vital and Non-Vital Nodes onto market behaviour of GME while controlling for network activity, such as the aggregated number of posts at a given time. %, which has been shown to be an important factor in modelling the influence of the network on markets \citep{Audrino2020,Gan2020,Jiao2020}. 

 %It is also remarkable that classical network measures, such as aggregated network activity, were unable to explain market movements as good as our designed change point variables. 
 \str{We find that the designed change point variables showed greater explanatory power for market movements than classical network measures, such as aggregated network activity.} 
 The vital users behavioural changes impact the market with a time lag of around 6 hours whereas the non-vital users impact the market immediately. Both the vital and non-vital users behavioural changes significantly relate to the integrated volatility even when controlling for aggregated network activity and established network measures proxying influential user impact. For the integrated volatility, the adjusted $R^2$ of a model without vs. with variables representing the behavioural change of individual users improves by 140\% (from 0.182 to 0.254). We also show that the change in behaviour is not explained by market dynamics, instead it is rooted in the activity taking place on the social network. 

%Analysing individual users influence and changes in the influence, is important for understanding how a network develops \citep{Aral2014, Li2011}. 

With our study, we contribute to a number of literatures. First, we provide empirical justification for theoretical models suggesting a link between social networks and financial markets \str{by quantifying the impact individual social network users have upon the markets} \citep{Pedersen2022}. Second, we contribute to the literature on link prediction in large online social networks by designing scalable node-specific link prediction models involving node behavioural and network variables \citep{bhattacharya2019coevolution, Lu2013, han2020importance}. Third, we contribute to the change point literature by proposing an algorithm for node-specific behavioural change detection and by studying the relationship between stock market movements and these changes \citep{chen2024monitoring, bai2023multiple}. Fourth, we contribute to the growing literature on meme stocks. \citep{Li2024,Aloosh2021, Costola2021,Klein2022,Umar2021}. Further, we provide the ISMB modelling approach which can be used for further understanding of \str{the impact individual social network users have upon the markets} beyond the effect of WSB on GME. By its design ISMB can be used not only to analyse meme stocks, but also meme coins, which emerge from similar situations like meme stocks. Hence, its applicability extends from stocks towards the cryptocurrency market. 

%Our paper has important managerial implications. First, our findings provide implications for Reddit and other online platforms. Understanding the relationship between online communications on social networks and stock markets helps platforms understand the importance of their role in the economy. The result that vital and non-vital users behaviour significantly predicts market movements, suggests that changes in stock market activity originate from users behavioural changes. This illustrates that focusing on individual users behaviour in addition to aggregated network activity is important. %Further the developed network model outperforms established metrics for measuring influential user impact by providing a framework to study a nodes behaviour over time. 
%Second, meme stocks are having a real impact on funds’ values; however, this influence is extremely unpredictable and volatile. Thus, for fund and portfolio management, a thorough understanding of network evolution and the underlying driving factors is crucial. Our model provides important insights on the relationship between social network discussions and stock prices. Third, a better understanding of the relationship between online social networks and financial markets is crucial for regulators. Our findings pave the way towards the construction of early warning signals for meme stock related market frenzies. 

The remainder of this paper is structured as follows. %Section \ref{sec:relwork} summarizes the relevant literature and presents the theoretical background. 
Section \ref{sec:data_desc} describes the WSB network data. Section \ref{sec:method} introduces our model and section \ref{sec:results} presents the results. Finally, section \ref{sec:conclusion} discusses the implications and conclusions of the study. In addition, we provide an Appendix \ref{appendix} presenting computational aspects and robustness checks.

%%%%%%%%%%%%% Old from former introduction

%In this paper we present a network model to study Reddit users posting behaviour as well as changes in their behaviour within a subreddit. The model rests on post and comment level data collected from Reddit as well as studies investigating the developments on Reddit. 
%The increasing importance of Reddit for markets and society likewise is underscored by studies surrounding its developments. 
%Research covers topics such as subreddit developments \citep{Thukral2018}, factors for influencers on subreddits \citep{Gianstefani2022}, network metrics \citep{Datta2017}, sentiment \citep{Umar2021}, financial literacy \citep{Klein2022}, and price impact \citep{Pedersen2022}. 
%It was studied how trends on subreddits contribute to its development \citep{Thukral2018}, the determining factors for influencers on subreddits \citep{Gianstefani2022} and how network metrics differ across subreddits \citep{Datta2017}. The sentiment expressed on the subreddit `Wallstreetbets' was widely studied \citep{Umar2021}, improvements to financial literacy \citep{Klein2022} and the price impact emerging from different trader types as by market beliefs \citep{Pedersen2022}. 




%Since the GameStop market frenzy ever new subreddits emerge in which users discuss investments, such as `stonks' and `superstonk' to name a selected few. Investors and the media equally recognise subreddits as a new driving force behind market movements which coined the term `meme stocks', describing assets being vividly discussed on social media. Owing to the hype around these stocks, a herding behaviour emerges which impacts asset prices, also reflected in the model presented by \textcite{Pedersen2022}. The increasing importance of Reddit for markets and society likewise motivates the need to study user posting behaviour on the platform, how they influence each other, and what drives changes in their behaviour over time. Understanding individual users behaviour enables one to identify users vital for the network, study their impact on other users, study and understand what drives the overall changes in the network which, as it is the case for meme stocks, influences the stock markets. 


%The resulting model has the advantage of having variables directly linked to the nodes' and networks' present and past behaviour, allowing to explain how they determine communication probability and what changed once the nodes posting probability changes. Contrary to many other studies, we only consider variables which are directly measurable and therefore do not have an estimation bias by themselves. 
%Many metrics considered in the financial network literature, such as sentiment and topic modelling, have an inherent estimation bias resulting from the chosen method. 
%We focus upon modelling and explaining the probability that a communication between two nodes takes place, a metric directly observable from the network, hence our model does not contain such estimation bias.  

%For each edge, the model is computationally highly efficient. 
%The model we are introducing is highly flexible as we derive a dedicated model for each edge connecting any two nodes. The considered parameters represent the past in and outgoing communication of a node, the communication of vital nodes it is associated with, and the activity taking place within the nodes all past and present posts. Since the overall activity of the network may also impact a nodes probability to communicate, we further embed present and past network activity as parameters to the model. Naturally one has to be awake and idle to communicate on a social network, hence we include hourly dummy variables to the model to control for respective effects. Lastly, since we target to model communication linked to financial markets, we include a dummy variable for the opening hours of the market. We model the probability of a communication taking place between two nodes with a logit model containing aforementioned variables. As not each of these variables is relevant for explaining the communication between every two nodes, we estimate the model with LASSO

%The most famous and probably best documented case of Reddit users impacting meme stocks, is the GME market frenzy with the Wallstreetbets subreddit at its core. Before and during the GME/WSB frenzy, the situation was very different than for any meme stock after GME. WSB was the main forum where the communication about GME took place. Since meme stocks were not a known concept yet, there was no major media coverage about it. After GME, new forums on Reddit were launched to discuss meme stocks and the media regularly reports about the topic. But for GME, it is established that WSB was the dominant platform for the investment decisions of retail traders. Hence it allows to study the effect social network users interactions have on markets, given the communication took place at a central location and the effects are free from mainstream media impact, as it was not covered yet. Hence, we will motivate and evaluate the developed network model on this subreddit to improve the understanding of the network microstructure, its change over time and the effects on the market. We obtain all the submits (posts and comments) on `Wallstreetbets' which mention `GameStop' or `GME' (the GameStop market ticker) in the title of the corresponding post. We obtain the data for the time period 01.06.2020 until 31.05.2021, whereas the market frenzy took place in January 2021. 
%It has been shown that network activity such as aggregated number of posts at a given time is an important factor in modeling network influence on markets \parencite{Audrino2020, Gan2020, Jiao2020}. While controlling for these and other effects, we study the relationship between the changepoints in the behaviour of individual vital and non-vital users and the market behaviour of GME. We analyse the effect of users change in behaviour on the integrated and jump volatility which we derive from hourly and 5 minute realized volatility, hourly and 5 minute trading volume as well as the corresponding returns. 
%The results show that the change in behaviour of vital and non-vital nodes strongly improves the predictability of said metrics. The adj. $R^2$ from a model without vs. with variables representing behavioural change of individual users, improves by $150\%$. During the market frenzy, these variables contribute towards reaching $0.45$ adj. $R^2$. Likewise for integrated variance, the changepoints of the vital/non-vital nodes excel in explaining market dynamics with the onset of the market frenzy. For jump volatility, even $0.8$ adj. $R^2$ is reached. Even better market predictability is reached when considering to build the models as by trading hour. It is further remarkable, that classical network measures such as AC and the network statistics were not able to explain the market movements that well. The explainability of the return dynamics from the changes of users behaviour remained high also after the frenzy, outperforming all other models. This shows that the social network had a longer lasting effect on the markets, going beyond the main market frenzy. 
%The results show that the integrated variance is significantly affected by vital users behavioural change with a time lag of 6 hours. Non-vital users significantly impact the integrated variance at various lags as well. As we found that the change in non-vital users behaviour is linked to vital users behavioural change, this translates into network effects impacting the markets. These results are robust when controlling for past integrated variance, aggregated network activity and network statistics such as mean distance, mean degree, number of active nodes and betweenness. It is also robust when accounting for hour of the day effects. Hence we follow that the regular volatility of the GME return (integrated variance) was impacted by the change in behaviour of the network users. Performing a similar analysis on the jump volatility, we find that the network statistics and network activity are a better predictor for the jump volatility than the change in users behaviour. Consequently the users of `Wallstreetbets' had a persistent impact on the volatility via the integrated variance whereas changes in the overall form of the network (network statistics, activity) impacted the spurious jump volatility instead. Likewise we find the return of the GME series and the change in trading volume being affected by the change in users behaviour.




